% NS-46 analytic proof draft. No numerical or cofinal sign claim.
\subsection*{Finite fine-lattice Grams without an infinite-lattice assumption}

\begin{lemma}[An elementary finite-arc Gram lower bound]
\label{lem:ns46-finite-arc}
Let $\phi\in L^2(\mathbb R)$ have unitary Fourier transform satisfying
$|\widehat\phi(\xi)|^2\ge c_0>0$ on $[-r,r]$, with $r>0$.
Let $m\ge2$, $\delta>0$ and $0<r\delta\le\pi/2$.
The Gram matrix of the $m$ translates
$\phi(\cdot-t_0-j\delta)$, $0\le j<m$, obeys
\begin{equation}\label{eq:ns46-gram-bound}
 \Gamma_{m,\delta}\succeq\gamma_{m,\delta}I,
 \qquad
 \gamma_{m,\delta}:=\frac{c_0r}{m}
       \left(\frac{r\delta}{4\pi e}\right)^{2(m-1)}.
\end{equation}
No positive lower bound for an infinite-lattice Gram is assumed.
The displayed bound is allowed to deteriorate with both $m$ and $\delta$.
\end{lemma}
\begin{proof}
For coefficients $c_j$ put $P(z)=\sum_{j=0}^{m-1}c_jz^j$ and
$\theta_0=r\delta$, $\ell=\theta_0/m$. The uncut physical norm is at least
\[
 \frac{c_0}{\delta}\int_{-\theta_0}^{\theta_0}
                      |P(e^{i\theta})|^2\,d\theta.
\]
The sign of the phase and the common shift do not affect this inequality.
Write $I_P$ for the integral. For each $0\le k<m$, the interval
\[
 B_k=[-\theta_0+2k\ell,-\theta_0+(2k+1)\ell]
\]
has length $\ell$, and lies inside the integration arc. Choose
$\theta_k\in B_k$ with
$|P(e^{i\theta_k})|\le\sqrt{I_P/\ell}$, which is possible by averaging.
For $j\ne k$ the angles satisfy
\[
 |\theta_j-\theta_k|\ge |j-k|\ell,
 \qquad |\theta_j-\theta_k|\le2\theta_0\le\pi.
\]
Thus, with $z_k=e^{i\theta_k}$,
$|z_j-z_k|\ge(2\ell/\pi)|j-k|$.

Lagrange interpolation at these distinct nodes is exact for $P$.
On the full unit circle its numerator factors have modulus at most two,
so
\begin{align*}
 |P(z)|
 &\le\sqrt{I_P/\ell}
   \sum_{k=0}^{m-1}
   \frac{2^{m-1}}{(2\ell/\pi)^{m-1}k!(m-1-k)!}\\
 &=\sqrt{I_P/\ell}\,
   \frac{(2\pi/\ell)^{m-1}}{(m-1)!}.
\end{align*}
Parseval on the full circle bounds $\sum_j|c_j|^2$ by the square
of this supremum. Therefore
\[
 I_P\ge\ell
 \left((m-1)!\left(\frac{\ell}{2\pi}\right)^{m-1}\right)^2
 \sum_j|c_j|^2.
\]
The elementary factorial bound
$(m-1)!\ge((m-1)/e)^{m-1}$ and $(m-1)/m\ge1/2$ give
\[
 (m-1)!\left(\frac{r\delta}{2\pi m}\right)^{m-1}
 \ge\left(\frac{r\delta}{4\pi e}\right)^{m-1}.
\]
Multiplying by $c_0/\delta$ proves the lemma, for complex coefficients
as well as real ones. The result uses only a finite polynomial and a
single fixed Fourier interval; it makes no assertion about the zeros
of the transform on any infinite arithmetic progression.
\end{proof}

\paragraph{Application data for the actual source.}
The Gaussian arithmetic radical in v1.35 is real, even, smooth and
nonnegative, and not zero, as made explicit in v1.50. Its normalized
version $\phi$ has $\widehat\phi(0)>0$. Continuity therefore supplies
fixed $r,c_0>0$ as in the lemma. Their existence is analytic; no numerical
value or certified effective threshold is asserted here.

For $a=\log\lambda$ sufficiently large, take
\[
 J_a=\left\lfloor\frac{\lambda^2}{40000a}\right\rfloor,
 \qquad m_a=2J_a+1,\qquad \delta_a=1/J_a,
 \qquad t_{a,j}=j\delta_a\quad(-J_a\le j\le J_a).
\]
All centers are in the fixed interval $[-1,1]$.
The lemma applies after a common shift to the indices $0,\ldots,m_a-1$.
For sufficiently large $a$,
\[
 m_a\le\frac{\lambda^2}{10000a},\qquad
 \log\left(\frac{4\pi e}{r\delta_a}\right)\le3a.
\]
Consequently its uncut Gram lower bound satisfies
\begin{equation}\label{eq:ns46-gram-asymptotic}
 \gamma_a\ge\frac{c_0r}{m_a}
       \exp\left(-\frac{6\lambda^2}{10000}\right).
\end{equation}
This is a lower bound for the smallest ordinary Gram eigenvalue,
not for a Weil eigenvalue. In particular the conditioning is explicitly
paid for, not silently assumed uniform.
